\documentclass[10pt]{article}
\usepackage[T1]{fontenc}
\usepackage[utf8]{inputenc}
\usepackage[french]{babel}
\usepackage{amsmath,amssymb,mathrsfs,mathpazo}
\usepackage{geometry}
\geometry{a4paper,margin=2cm}
\usepackage{graphicx,array,tabularx,tabularray,colortbl,multirow,multicol,booktabs,enumitem,pifont,ulem,fancybox,diagbox,cancel}
\usepackage{pstricks,pstricks-add,pst-plot,pst-tree,pst-node,pst-text,pst-func,pst-math,pst-bspline,pst-3dplot,pst-eucl,pst-solides3d}
\let\pscalculate\undefined
\usepackage{pst-bezier}
\usepackage{tikz,pgfplots,tkz-tab}
\usetikzlibrary{arrows,arrows.meta,patterns,calc,positioning,trees,decorations.markings,shapes.geometric,shapes.misc,shapes.arrows,intersections,decorations.pathmorphing,decorations.pathreplacing,backgrounds,shadows,babel}
\pgfplotsset{compat=1.18}
\usepgfplotslibrary{fillbetween}
\usepackage[np]{numprint}
\usepackage{listings,makecell,pas-tableur,twoopt}
\let\barre\undefined
\usepackage[pstricks]{tablvar}
\definecolor{darkgreen}{RGB}{0,100,0}
\providecommand{\red}{\color{red}}
\providecommand{\blue}{\color{blue}}
\providecommand{\green}{\color{green}}
\providecommand{\black}{\color{black}}
\providecommand{\white}{\color{white}}
\providecommand{\gray}{\color{gray}}
\providecommand{\cyan}{\color{cyan}}
\providecommand{\magenta}{\color{magenta}}
\providecommand{\yellow}{\color{yellow}}
\providecommand{\violet}{\color{violet}}
\providecommand{\orange}{\color{orange}}
\providecommand{\euro}{\texteuro}
\providecommand{\EUR}{\texteuro}
\providecommand{\ieme}{\textsuperscript{e}}
\providecommand{\degres}{\ensuremath{^\circ}}
\providecommand{\degree}{\ensuremath{^\circ}}
\providecommand{\vv}[1]{\overrightarrow{#1}}
\providecommand{\Coord}[1]{#1}
\providecommand{\decosix}{}
\providecommand{\starredbullet}{\ensuremath{\bullet}}
\pagestyle{empty}
\setlength{\parindent}{0pt}
\usepackage{adjustbox}\DeclareUnicodeCharacter{2236}{:}\DeclareUnicodeCharacter{2190}{\ensuremath{\leftarrow}}
\providecommand{\R}{\mathbb{R}}
\providecommand{\N}{\mathbb{N}}
\providecommand{\Z}{\mathbb{Z}}
\providecommand{\Q}{\mathbb{Q}}
\providecommand{\D}{\mathbb{D}}
\providecommand{\C}{\mathbb{C}}
\providecommand{\e}{{\mathrm{e}}}
\providecommand{\vect}[1]{\overrightarrow{#1}}
\providecommand{\vectt}[1]{\overrightarrow{\mathrm{#1}}}
\providecommand{\vvt}[1]{\overrightarrow{\mathrm{#1}}}
\providecommand{\np}[2][]{{#2~#1}}
\providecommand{\vv}[1]{{\overrightarrow{#1}}}
\renewcommand{\d}{{\mathrm{d}}}
\providecommand{\degres}{{^\circ}}
\providecommand{\degre}{{^\circ}}
\providecommand{\blue}{\color{blue}}
\providecommand{\red}{\color{red}}
\providecommand{\enskip}{\quad}
\providecommand{\llbracket}{\lbrack\!\lbrack}
\providecommand{\rrbracket}{\rbrack\!\rbrack}
\providecommand{\euro}{{\text{€}}}
\providecommand{\no}{{\text{n°}}}
\providecommand{\interval}[2]{{[#1;#2]}}
\providecommand{\intervFO}[2]{{[#1;#2[}}
\providecommand{\dotfill}{\ldots}

\providecommand{\ds}{\displaystyle}
\providecommand{\textding}[1]{#1}
\providecommand{\Oij}{\(\mathrm{O};\vec{\imath},\vec{\jmath}\)}
\providecommand{\Oijk}{\(\mathrm{O};\vec{\imath},\vec{\jmath},\vec{k}\)}
\providecommand{\Ouv}{\(\mathrm{O};\vec{u},\vec{v}\)}
\providecommand{\intervFF}[2]{\ensuremath{[#1;#2]}}
\renewcommand{\euro}{\ifmmode\text{\texteuro}\else\texteuro\fi}
\DeclareUnicodeCharacter{00B0}{\ifmmode^{\circ}\else\textdegree\fi}
\DeclareUnicodeCharacter{20AC}{\ifmmode\text{\texteuro}\else\texteuro\fi}
% Énoncés : Éducation nationale. Transcription et corrigés : APMEP, Denis Vergès et contributeurs. Découpage et classement : CoopMaths, Sébastien Lozano, Éric Elter et contributeurs. Adaptation de la présentation : Maths Explorer.
% CC BY-SA : https://coopmaths.fr/www/about/
% https://coopmaths.fr/alea/static/e3c/2020/tex/e3c_2020_00_sujet40_4_cor.tex
\begin{document}
\noindent\textit{APMEP / CoopMaths. CC BY-SA. Présentation adaptée par Maths Explorer.}\par\medskip
\subsection*{1.}

\(f(x) = 4x^3 - 48x^2 + 144x\).

\paragraph{a.} La fonction polynôme \( f \) est dérivable sur \( \mathbb{R} \) et sur cet intervalle :
\[
f'(x) = 12x^2 - 96x + 144 = 12x^2 - 12 \times 8x + 12 \times 12 = 12(x^2 - 8x + 12).
\]

\paragraph{b.} Pour le trinôme \( x^2 - 8x + 12 \), on a :
\[
\Delta = 64 - 48 = 16 = 4^2.
\]

Il y a donc deux racines :
\[
x_1 = \dfrac{8 + 4}{2} = 6 \quad \text{et} \quad x_2 = \dfrac{8 - 4}{2} = 2.
\]

On sait que le trinôme est positif sauf sur l'intervalle \( ]2\,;\,6[ \) où \( f'(x) < 0 \).

La fonction est donc croissante sauf sur l'intervalle \( ]2\,;\,6[ \).

Avec \( f(2) = 128 \) et \( f(6) = 0 \), on a le tableau de variations suivant :
\begin{center}
\begin{tikzpicture}
\tkzTabInit[lgt=3, espcl=3.5]{$x$ / 1, {Signe de $f'(x)$} / 1, {$f$} / 2}{${-\infty}$, ${2}$, ${6}$, ${+\infty}$}
\tkzTabLine{,+,0,-,0,+,}
\tkzTabVar{-/{$$},+/{$128$},-/{$0$},+/{$$}}{/}
\end{tikzpicture}
\end{center}

\subsection*{2.}

\paragraph{a.} La boîte a une base carrée de \( 12 - 2 \times 1 = 10 \) de côté et une hauteur de 1. Son volume est donc égal à :
\[
V(1) = 1 \times 10^2 = 100 \, \text{cm}^3.
\]

\paragraph{b.} En général, pour \( 0 \leqslant x \leqslant 6 \), la boîte a une base carrée de \( 12 - 2x \) de côté et une hauteur de \( x \). Son volume est donc égal à :
\[
V(x) = (12 - 2x)^2 \times x = x \times (144 + 4x^2 - 48x) = 4x^3 - 48x^2 + 144x = f(x).
\]

\paragraph{c.} On a vu que \( f \) a un maximum : \( f(2) = 128 \).

Le volume maximal \( 128 \, \text{cm}^3 \) est obtenu pour \( x = 2 \, \text{cm} \).

\end{document}
